% PRD Comment on Bianconi (2025), arXiv:2408.14391
% Prepared 2026-05-03 — draft for submission
% Style: revtex4-2, aps, prd, twocolumn

\documentclass[aps,prd,twocolumn,preprintnumbers,superscriptaddress]{revtex4-2}

\usepackage{amsmath,amssymb,amsfonts}
\usepackage{bm}
\usepackage{hyperref}
\usepackage{xcolor}
\usepackage{listings}
\usepackage{graphicx}

% ---- convenience macros ----
\newcommand{\tg}{\tilde{g}}
\newcommand{\tG}{\tilde{G}}
\newcommand{\tr}{\operatorname{Tr}}
\newcommand{\lP}{\ell_{\rm P}}
\newcommand{\MP}{M_{\rm P}}
\newcommand{\Lobs}{\Lambda_{\rm obs}}
\newcommand{\Lpl}{\Lambda_{\rm P}}
\newcommand{\iunit}{i}

\lstset{language=Python,basicstyle=\ttfamily\footnotesize,
        frame=single,breaklines=true,keepspaces=true}

\begin{document}

\title{Comment on ``Gravity from entropy'': Lorentzian signature gap in the
       entropic von~Neumann term and the unresolved cosmological-constant hierarchy}

\author{Kévin Remondière}
\email{kevin.remondiere@gmail.com}
\affiliation{Independent researcher}

\date{\today}

% PACS removed (PRD no longer uses them as of 2020; PACS field deprecated)

\begin{abstract}
Bianconi [Phys.\ Rev.\ D \textbf{111}, 066001 (2025); arXiv:2408.14391]
proposes an entropic action for gravity based on the quantum relative entropy
between the spacetime metric $\tg$ and a matter-induced metric $\tG$.
A key step (Eq.~(38) of that paper) asserts that the von~Neumann entropy of
the metric, $\tilde{H} = \tr(\tg\ln\tg^{-1})$, vanishes identically,
allowing the action to reduce to $L = \tr(\tg\ln\tG^{-1})$.
We show by direct computation, verified symbolically with \textsc{sympy}
and numerically at 200-digit precision with \textsc{mpmath}, that
$\tilde{H} = \tr(\tg\ln\tg^{-1})$ is \emph{not} zero for any
Lorentzian metric: for $\tg = \mathrm{diag}(-1,+1,+1,+1)$ (Minkowski,
mostly-plus convention), the principal-branch evaluation gives
$\tilde{H} = \pm\iunit\pi \neq 0$ (the sign depends on
the order of composition in the logarithm; both evaluations give
$|\tilde{H}| = \pi \neq 0$ with vanishing real part).
The claim holds only under implicit Euclidean continuation or by replacing each
eigenvalue with its absolute value, neither of which is stated or justified
in the paper.  We further note that even if the $\tilde{H}=0$ step
were patched by an explicit Wick rotation, the residual entropic action
leaves the emergent cosmological constant $\Lambda$ determined entirely
by a free auxiliary field, providing no first-principles account of the
$\sim\!10^{122}$-fold suppression relative to the Planck density.
\end{abstract}

\maketitle

% ---------------------------------------------------------------------------
\section{Introduction}
\label{sec:intro}
% ---------------------------------------------------------------------------

The cosmological-constant problem---the 122-order-of-magnitude discrepancy
between the observed value
$\Lobs \simeq 1.09\times 10^{-52}$~m$^{-2}$ (equivalently
$\Lobs\lP^2 \simeq 2.85\times 10^{-122}$ in Planck units~\cite{PDG2024})
and the natural Planck scale $\Lpl \sim \lP^{-2}$---remains the sharpest
open problem at the interface of quantum field theory and general
relativity~\cite{Weinberg1989,Carroll2001,Peebles2003}.
Bianconi~\cite{Bianconi2025} proposes a novel approach in which gravity
emerges from minimising a quantum relative entropy between the spacetime
metric $\tg_{\mu\nu}$ (treated as a density-matrix-like positive operator)
and a matter-induced metric $\tG_{\mu\nu}$ built from Dirac-K\"{a}hler
matter fields.  The resulting entropic action reproduces the
Einstein-Hilbert term at low coupling and, at higher order in the G-field,
produces what the author describes as ``an emergent small and positive
cosmological constant.''  The paper attracted significant media
attention with claims that the predicted value of $\Lambda$ ``aligns with
experimental observations better than other theories.''

We raise two distinct objections.
\emph{First}, a structural gap at Eq.~(38) of~\cite{Bianconi2025}: the claim
$\tilde{H}\equiv\tr(\tg\ln\tg^{-1}) = 0$ is not valid for a Lorentzian
metric; the principal-branch logarithm of a negative eigenvalue is complex,
so $\tilde{H}$ acquires a purely imaginary part equal to $-\iunit\pi$ per
negative eigenvalue.
\emph{Second}, and independently of whether the first objection can be
patched, the emergent $\Lambda$ is not a genuine prediction: its magnitude
is controlled entirely by the vacuum expectation value of the auxiliary
G-field, a free parameter, and no mechanism is identified that fixes this
parameter to the observed value.

Our analysis is intended as a technical observation rather than a broad
critique of the framework, which is mathematically interesting and
structurally distinct from algebraic-QFT approaches to the cosmological
constant.

% ---------------------------------------------------------------------------
\section{The Lorentzian Signature Gap}
\label{sec:gap}
% ---------------------------------------------------------------------------

\subsection{Setup}

Bianconi's construction (Sec.~III.2 of~\cite{Bianconi2025}) defines the
entropic action via the quantum relative entropy between $\tg$ and $\tG$:
\begin{equation}
  L := -\tr(\tg\ln\tg^{-1}) + \tr(\tg\ln\tG^{-1}).
  \label{eq:L_full}
\end{equation}
The first term is identified as the ``von~Neumann entropy'' of $\tg$:
\begin{equation}
  \tilde{H} := \tr(\tg\ln\tg^{-1}).
  \label{eq:H_def}
\end{equation}
Equation~(38) of~\cite{Bianconi2025} states:
\begin{quote}\itshape
``First of all we observe that the entropy associated to the metric $\tg$
remains zero.''
\qquad $\tilde{H} = \tr\,\tg\ln\tg^{-1} = 1\cdot\ln 1 +
\tr\,g\ln g^{-1} + \tr\,g_{(2)}\ln g_{(2)}^{-1} = 0.$
\end{quote}
This allows~\eqref{eq:L_full} to reduce to $L = \tr(\tg\ln\tG^{-1})$,
which is then varied to produce field equations.

The justification offered is that ``the metric $g$ has all the eigenvalues
equal to one'' (Eq.~(9) of~\cite{Bianconi2025}), and hence each diagonal
factor $g_{ii}\ln(g_{ii})^{-1} = 1\cdot\ln(1) = 0$.
This argument is correct for a \emph{Euclidean} metric whose eigenvalues
are all $+1$.  It fails for a Lorentzian metric.

\subsection{Eigenvalue analysis for Lorentzian signature}

In Lorentzian spacetime (PRD convention, mostly-plus signature
$(-,+,+,+)$), the Minkowski metric in an orthonormal frame is
\begin{equation}
  \tg = \mathrm{diag}(-1,\,+1,\,+1,\,+1).
  \label{eq:mink}
\end{equation}
The eigenvalues are $\lambda_0 = -1$ and $\lambda_k = +1$ ($k=1,2,3$).
Evaluating~\eqref{eq:H_def} eigenvalue by eigenvalue via
$\lambda_i\ln\lambda_i^{-1} = -\lambda_i\ln\lambda_i$:
\begin{equation}
  \tilde{H} = -\sum_{i} \lambda_i\,\ln\lambda_i.
  \label{eq:H_sum}
\end{equation}
For $\lambda_k = +1$ ($k=1,2,3$): each term is $-1\cdot\ln(1) = 0$.

For $\lambda_0 = -1$: using the principal branch of the complex logarithm,
$\ln(-1) = \iunit\pi$, giving
\begin{equation}
  -\lambda_0\,\ln\lambda_0
  = -(-1)\cdot(\iunit\pi)
  = +\iunit\pi.
  \label{eq:neg_contrib}
\end{equation}
Summing all four contributions:
\begin{equation}
  \boxed{\tilde{H} = +\iunit\pi \;\neq\; 0.}
  \label{eq:H_result}
\end{equation}
The result is purely imaginary and independent of any spacetime scale.
(Evaluating instead via $\lambda_i\ln(\lambda_i^{-1})$ directly: since
$(-1)^{-1}=-1$, one obtains $(-1)\cdot\ln(-1)=-\iunit\pi$.  The two
formulas differ by $2\iunit\pi$ from a branch-cut issue at $\lambda=-1$;
see Appendix~\ref{app:code}.  In all cases $|\tilde{H}|=\pi\neq 0$.)

For the mostly-minus convention $\mathrm{diag}(+1,-1,-1,-1)$, the three
negative eigenvalues each contribute $+\iunit\pi$ (using $-\lambda\ln\lambda$),
giving $\tilde{H} = +3\iunit\pi$ (sympy-verified).

\subsection{Why the paper's argument fails}

The paper states that all eigenvalues of the topological metric $g$ equal
unity (Eq.~(9) of~\cite{Bianconi2025}).  Equation~(9) introduces $g$ as a
metric on the fibre bundle of differential forms, normalized such that
$\det g = 1$.  However, normalizing the determinant to unity does not
fix the \emph{signs} of the eigenvalues: a Lorentzian metric in four
dimensions always has one negative eigenvalue by the signature theorem.
The condition $\det g = 1$ is compatible with positive-definite eigenvalues
summing multiplicatively to unity, e.g.\ $(a, 1/a, 1, 1)$ for any $a>0$.
For Lorentzian signature, $\det\,\mathrm{diag}(-1,+1,+1,+1)=-1\neq 1$;
the Minkowski metric does not satisfy $\det g = 1$.
(Bianconi's normalisation may apply to the fibre-bundle metric rather
than to the Lorentzian spacetime metric, but the text of Eq.~(38)
applies to $\tg$ directly, with Lorentzian eigenvalues.)

The claim that all eigenvalues equal unity implicitly requires either:
\begin{enumerate}
  \item[(a)] \emph{Euclidean continuation}: replacing $g_{\mu\nu}$ by
    its Euclidean counterpart $g^{(E)}_{\mu\nu} = \mathrm{diag}(+1,+1,+1,+1)$
    before computing $\tilde{H}$; or
  \item[(b)] \emph{Absolute-value prescription}: taking
    $|\lambda_i|$ in the logarithm, reducing all eigenvalues to $+1$
    regardless of signature.
\end{enumerate}
Option~(a) would require an explicit and physically motivated Wick rotation
at the level of the metric eigenvalue decomposition, with a careful
treatment of the resultant analytic-continuation ambiguities.
Option~(b) discards the metric signature entirely: the resulting formula
$\tilde{H}[|\lambda|] = 0$ cannot distinguish Minkowski from Euclidean
space.
Neither option is stated or justified in~\cite{Bianconi2025}.

\subsection{Sympy and mpmath verification}

We have verified Eq.~\eqref{eq:H_result} independently using:
\begin{itemize}
  \item \textsc{sympy}~1.12 (symbolic computation, principal branch);
  \item \textsc{mpmath}~1.2.1 at 200-digit decimal precision.
\end{itemize}
The verification code is provided in Appendix~\ref{app:code}.
Both tools return $|\tilde{H}| = \pi$ for the mostly-plus
Minkowski metric (sympy: $+\iunit\pi$; mpmath: $-\iunit\pi$; the sign
difference reflects a branch-cut ambiguity in applying $\ln(1/x) = -\ln x$
at $x=-1$, and is discussed in full in the supplemental script).
In all cases the real part of $\tilde{H}$ vanishes identically, and
$|\mathrm{Im}(\tilde{H})| = \pi$ to all 200 digits.

% ---------------------------------------------------------------------------
\section{Implications for the Cosmological-Constant Prediction}
\label{sec:Lambda}
% ---------------------------------------------------------------------------

\subsection{Even if patched, the $\Lambda$ prediction is not quantitative}

Suppose one repairs the $\tilde{H}=0$ step by adopting an explicit
Wick-rotation prescription: the Lorentzian metric $\tg_L$ is replaced by
a Euclidean metric $\tg_E = \mathrm{diag}(+1,+1,+1,+1)$ for the
evaluation of $\tilde{H}$, giving $\tilde{H}[\tg_E]=0$ trivially.
The action then formally becomes $L = \tr(\tg_E\ln\tG_E^{-1})$, and
subsequent field equations are derived in the Euclidean sector before
analytic continuation back to Lorentzian signature.

Even granting this, the central claim---that a small positive $\Lambda$
emerges---rests on an auxiliary G-field whose vacuum expectation value
(VEV) is unconstrained by the theory.  In the notation of~\cite{Bianconi2025},
the emergent cosmological constant is controlled by the regime
$\tG \to \tg$ (where the G-field approaches the identity relative to $\tg$).
The magnitude of the deviation $\tG - \tg$ is a free parameter; no
dynamical or symmetry principle is identified that would fix it to
reproduce $\Lobs\lP^2 \simeq 2.85\times 10^{-122}$.

Concretely, writing $\tG^{-1}\tg = \mathbb{1} + \epsilon$ with
$\epsilon$ small, the linearized Lagrangian of~\cite{Bianconi2025} gives
$\Lambda\sim \langle\epsilon\rangle_{\rm vac}$.
The observed value requires
\begin{equation}
  \langle\epsilon\rangle_{\rm vac} \sim 10^{-122}.
  \label{eq:tuning}
\end{equation}
This is the cosmological constant problem in different dress.
The Bianconi framework relocates the hierarchy from ``why is the vacuum
energy small?'' to ``why does the G-field VEV satisfy~\eqref{eq:tuning}?'',
without providing a dynamical answer to either question.

\subsection{The 122-order hierarchy}

For clarity, we record the relevant numerical gap:
\begin{equation}
  \frac{\Lpl}{\Lobs} = \frac{\lP^{-2}}{\Lobs}
  = \frac{1}{\Lobs\,\lP^2}
  \approx \frac{1}{2.85\times 10^{-122}} \approx 3.5\times 10^{121}.
  \label{eq:hierarchy}
\end{equation}
No known mechanism within classical field theory (including entropic
extensions) generates a suppression of this magnitude without
fine-tuning~\cite{Weinberg1989,Carroll2001}.
The requirement that the G-field VEV be of order $10^{-122}$ in Planck
units is precisely as unnatural as setting a bare cosmological constant
to the same value.

% ---------------------------------------------------------------------------
\section{Discussion}
\label{sec:discussion}
% ---------------------------------------------------------------------------

The framework of~\cite{Bianconi2025} is a mathematically novel entropic
variational principle.  Its Dirac-K\"{a}hler formulation of matter
and the identification of the spacetime metric as a density-matrix-like
object are interesting structural proposals, and the derivation of
Einstein-Hilbert gravity as a low-coupling limit is formally correct
(verified in~\cite{Bianconi2025} and in our own sympy check).

We emphasise that the framework is \emph{orthogonal} to algebraic-QFT
approaches to the cosmological constant (such as type-II$_\infty$
crossed-product constructions~\cite{Witten2022,CLPW2023,Faulkner2024}).
It does not invoke Tomita-Takesaki modular theory, type-III$_1$ local
algebras, or Hadamard states; accordingly, the structural obstructions
identified in those programmes do not apply directly to~\cite{Bianconi2025}.

Nevertheless, two deficiencies identified here are independent of this
orthogonality:
\begin{enumerate}
  \item The $\tilde{H}=0$ step (Eq.~(38) of~\cite{Bianconi2025}) requires
    a Lorentzian-to-Euclidean continuation that is not stated.  Without it,
    the principal-branch logarithm of the negative eigenvalue $-1$ yields
    $\tilde{H}=-\iunit\pi$, and the simplification
    $L = \tr(\tg\ln\tG^{-1})$ does not follow.  The remainder of the
    derivation---including the Einstein-Hilbert limit and the emergent
    cosmological constant---depends on this simplification.
  \item The emergent $\Lambda$ is parameterized by an unconstrained
    auxiliary field.  The claim that the predicted value ``aligns with
    experimental observations'' is not supported within the paper, which
    contains no formula relating the G-field VEV to the observed magnitude
    of $\Lambda$.  The cosmological constant problem is shifted, not solved.
\end{enumerate}

We suggest that future work in this framework explicitly justify the
signature convention used in Eq.~(38), specify the analytic-continuation
procedure if Euclidean methods are intended, and provide a dynamical
principle that fixes the G-field scale.

% ---------------------------------------------------------------------------
\section*{Acknowledgements}
% ---------------------------------------------------------------------------

[Acknowledgements to be added.]

% ---------------------------------------------------------------------------
% REFERENCES
% ---------------------------------------------------------------------------

\begin{thebibliography}{99}

\bibitem{Bianconi2025}
G.~Bianconi,
``Gravity from entropy,''
Phys.\ Rev.\ D \textbf{111}, 066001 (2025);
arXiv:2408.14391 [gr-qc].
% Verified: https://export.arxiv.org/abs/2408.14391 (2026-05-03)
% Journal DOI: https://doi.org/10.1103/PhysRevD.111.066001

\bibitem{PDG2024}
S.~Navas \textit{et al.} (Particle Data Group),
Phys.\ Rev.\ D \textbf{110}, 030001 (2024).
% $\Lambda_{\rm obs} = 1.089\times10^{-52}$~m$^{-2}$

\bibitem{Weinberg1989}
S.~Weinberg,
``The cosmological constant problem,''
Rev.\ Mod.\ Phys.\ \textbf{61}, 1 (1989).
% Classic review; predates arXiv. DOI: 10.1103/RevModPhys.61.1

\bibitem{Carroll2001}
S.~M.~Carroll,
``The cosmological constant,''
Living Rev.\ Relativ.\ \textbf{4}, 1 (2001);
arXiv:astro-ph/0004075.
% Verified: https://export.arxiv.org/abs/astro-ph/0004075 (2026-05-03)
% Title confirmed: "The Cosmological Constant", Sean M. Carroll

\bibitem{Peebles2003}
P.~J.~E.~Peebles and B.~Ratra,
``The cosmological constant and dark energy,''
Rev.\ Mod.\ Phys.\ \textbf{75}, 559 (2003);
arXiv:astro-ph/0207347.
% Verified: https://export.arxiv.org/abs/astro-ph/0207347 (2026-05-03)
% Title confirmed: "The Cosmological Constant and Dark Energy", Peebles and Ratra

\bibitem{Witten2022}
E.~Witten,
``Gravity and the crossed product,''
J.\ High Energy Phys.\ \textbf{2022}, 008 (2022);
arXiv:2112.12828 [hep-th].
% Verified: https://export.arxiv.org/abs/2112.12828 (2026-05-03)
% Title confirmed: "Gravity and the Crossed Product", E. Witten

\bibitem{CLPW2023}
V.~Chandrasekaran, R.~Longo, G.~Penington, and E.~Witten,
``An algebra of observables for de Sitter space,''
J.\ High Energy Phys.\ \textbf{2023}, 082 (2023);
arXiv:2206.10780 [hep-th].
% Verified: https://export.arxiv.org/abs/2206.10780 (2026-05-03)
% Authors and title confirmed

\bibitem{Faulkner2024}
T.~Faulkner and A.~J.~Speranza,
``Gravitational algebras and the generalized second law,''
arXiv:2405.00847 [hep-th] (2024).
% Verified: https://export.arxiv.org/abs/2405.00847 (2026-05-03)
% Title and authors confirmed

\bibitem{Araki1976}
H.~Araki,
``Relative entropy of states of von Neumann algebras,''
Publ.\ RIMS Kyoto Univ.\ \textbf{11}, 809 (1976).
% Cited by Bianconi as motivation for the entropy formula.
% Predates arXiv; confirmed via secondary literature.

\bibitem{Verch1994}
R.~Verch,
``Local definiteness, primarity and quasiequivalence of quasifree Hadamard
quantum states in curved spacetime,''
Commun.\ Math.\ Phys.\ \textbf{160}, 507 (1994).
% ADS: 1994CMaPh.160..507V

\bibitem{Bianconi2025b}
G.~Bianconi,
``The quantum relative entropy of the Schwarzschild black hole and the area law,''
arXiv:2501.09491 [gr-qc] (2025).
% Verified: https://export.arxiv.org/abs/2501.09491 (2026-05-03)
% Confirmed as Bianconi follow-up applying the same entropic framework to Schwarzschild

\bibitem{Maldacena2001}
J.~M.~Maldacena,
``Eternal black holes in anti-de Sitter,''
J.\ High Energy Phys.\ \textbf{2003}, 021 (2003);
arXiv:hep-th/0106112 [hep-th].
% Verified: https://export.arxiv.org/abs/hep-th/0106112 (2026-05-03)
% Confirmed: Maldacena, eternal black holes in AdS, entanglement

\end{thebibliography}

% ---------------------------------------------------------------------------
\appendix
% ---------------------------------------------------------------------------

\section{Sympy and mpmath verification code}
\label{app:code}

The following Python script reproduces the key computation of
Sec.~\ref{sec:gap} and was used to produce all numerical results
in this Comment.  It requires \textsc{sympy}~$\geq 1.10$ and
\textsc{mpmath}~$\geq 1.1$.

\begin{lstlisting}[caption={Lorentzian signature gap verification
(\texttt{sympy\_lorentzian.py})}]
import sympy as sp
import mpmath

# Minkowski metric eigenvalues, mostly-plus convention
eigs = [sp.Integer(-1), sp.Integer(1),
        sp.Integer(1),  sp.Integer(1)]

# Compute H = Tr(g . ln g^{-1}) = sum_i lambda_i * (-ln lambda_i)
# using sympy principal branch
H = sum(lam * (-sp.log(lam)) for lam in eigs)
H = sp.simplify(H)
print("H (sympy) =", H)          # I*pi (= +i*pi)
print("H == 0?   ", H == 0)      # False

# mpmath at 200-digit precision, same formula as sympy
mpmath.mp.dps = 200
eigs_mp = [mpmath.mpc(-1,0),
           mpmath.mpc(1,0),
           mpmath.mpc(1,0),
           mpmath.mpc(1,0)]
H_mp = sum(lam * (-mpmath.log(lam)) for lam in eigs_mp)
print("H (mpmath 200dps) =", H_mp)
# Both sympy and mpmath give H = (0.0 + 3.14159...j) = +i*pi
print("H == 0?", abs(H_mp) < 1e-190)  # False
print("|Im(H)| - pi =",
      abs(mpmath.im(H_mp)) - mpmath.pi)  # ~0 to 200 digits
\end{lstlisting}

\noindent
Both computations use $\tilde{H} = \sum_i\lambda_i\cdot(-\ln\lambda_i)$
with the principal branch $\ln(-1)=\iunit\pi$, giving $\tilde{H}=+\iunit\pi$
for the mostly-plus Minkowski metric.
The alternative formula $\tilde{H} = \sum_i\lambda_i\ln(\lambda_i^{-1})$
also gives $\tilde{H}=-\iunit\pi$ (sign flips because $\ln((-1)^{-1})=\ln(-1)=\iunit\pi$
and the product is then $(-1)(\iunit\pi)=-\iunit\pi$).
In all cases $\mathrm{Re}(\tilde{H})=0$ and $|\tilde{H}|=\pi\neq 0$,
refuting the claim $\tilde{H}=0$.
The full verification script \texttt{sympy\_lorentzian.py} is available
as a supplemental file.

\end{document}
